\uuid{Khcl}
\titre{Primitives usuelles}
\auteur{hanani ; mzari}
\organisation{exo7}
\licence{CC BY-NC-SA 4.0}
\datecreate{2020-01-20}
\niveau{L1}
\module{Analyse}
\chapitre{Calcul d'intégrales}
\sousChapitre{Primitives diverses}
\difficulte{2}
\theme{Primitives, primitives usuelles}
\exo7qcmid{lille-2.primitives.022}

\contenu{
\begin{question}
\qclassification{127.11}
\qtags{motcle=primitives usuelles}
Parmi les égalités suivantes, cocher celles qui sont vraies :
\begin{answers}
\good{$\displaystyle \int x\cos(1+x^2)\mathrm{d}x=\frac{1}{2}\sin(1+x^2)+k$, où $k\in \Rr$.}
\bad{$\displaystyle \int \mathrm{e}^x\sin\left(2+\mathrm{e}^x\right)\mathrm{d}x=\int \mathrm{e}^x\mathrm{d}x\times \int\sin\left(2+\mathrm{e}^x\right)\mathrm{d}x$.}
\bad{$\displaystyle \int x\cos(1+x^2)\mathrm{d}x=\frac{x^2}{2}\sin(1+x^2)+k$, où $k\in \Rr$.}
\good{$\displaystyle \int \mathrm{e}^x\sin\left(2+\mathrm{e}^x\right)\mathrm{d}x=-\cos\left(2+\mathrm{e}^x\right)+k$, où $k\in \Rr$.}
\end{answers}
\begin{explanations}
Avec $u=1+x^2$, on a : $\mathrm{d}u=2x\mathrm{d}x$ et
$$\int x\cos(1+x^2)\mathrm{d}x=\frac{1}{2}\int \cos (u)\mathrm{d}u=\frac{1}{2}\times \sin (u)+k=\frac{1}{2}\sin(1+x^2)+k,\; k\in \Rr.$$
De même, avec $u=2+\mathrm{e}^x$, on a : $\mathrm{d}u=\mathrm{e}^x\mathrm{d}x$ et
$$\int \mathrm{e}^x\sin\left(2+\mathrm{e}^x\right)\mathrm{d}x=\int \sin (u)\mathrm{d}u=-\cos (u)+k=-\cos\left(2+\mathrm{e}^x\right)+k,\; k\in \Rr.$$
\end{explanations}
\end{question}
}
